%!TeX root=MemoriaTFG.tex

\chapter*{Resum}
\addcontentsline{toc}{chapter}{Resum}

El Truc és un joc de cartes tradicional d'arrel catalana i valenciana, molt arrelat a
les Illes Balears, que es juga per parelles i combina informació oculta, dues apostes
independents (l'envit i el truc) i el farol. Aquesta barreja el converteix en un problema
d'informació parcial (formalment, un procés de decisió de Markov parcialment observable)
especialment exigent per a la intel·ligència artificial i, alhora, en un patrimoni cultural
amenaçat per la dificultat
creixent de reunir-hi jugadors. Aquest treball desenvolupa i avalua agents d'aprenentatge
per reforç profund capaços de jugar al Truc de manera competitiva, partint de models amb pesos aleatoris i sense
coneixement específic del joc injectat prèviament, en la variant d'un contra un sense senyes.

El treball s'articula en tres parts. Primer es construeix un simulador complet del joc amb
arquitectura Model-Vista-Controlador, jugable per humans i exposat com a entorn estàndard compatible amb
les biblioteques d'aprenentatge per reforç, amb una observació de 240~dimensions que combina
un tensor de cartes i un vector de context. Després s'aplica una metodologia experimental
incremental de sis fases, cadascuna guiada per una hipòtesi que condiciona la següent.
Aquesta inclou una
comparativa homogènia d'algorismes (\acs{DQN}, \acs{PPO} i \acs{NFSP}), \emph{curriculum
learning} de mans a partides senceres, un extractor de característiques convolucional amb
preentrenament supervisat, memòria recurrent amb un conjunt d'adversaris, \emph{self-play}
mixt i, finalment, \acs{NFSP} complet per aproximar-se a l'equilibri de Nash.

Els resultats indiquen que el coll d'ampolla inicial no és el còmput sinó el senyal
d'aprenentatge: entrenar primer sobre mans individuals incrementa el win-rate ponderat de
\acs{PPO} del 3{,}8\% al 84{,}8\%, i un cos convolucional preentrenat i congelat arriba al
91{,}2\%. El \emph{self-play} mixt millora la robustesa davant estils d'oponent diversos
sense sacrificar rendiment, i el \acs{NFSP} assoleix un punt d'equilibri de Nash genuí. El
model desplegat bat totes les variants de l'oponent heurístic i empata contra si mateix. Ara
bé, l'anàlisi qualitativa mostra que la proximitat estadística a l'equilibri no equival a un
rival creïble i inexplotable per a una persona: la gestió de l'envit és sòlida, però el farol
al truc és sistemàtic i, per tant, llegible. L'extensió a la variant cooperativa de dos
contra dos amb senyes queda com a principal línia de treball futur.

\bigskip
\noindent\textbf{Paraules clau:} \emph{reinforcement learning, truc, jocs de cartes, DQN, PPO, NFSP, informació parcial, curriculum learning, self-play.}
